\begin{proposition}[确定性价值函数的凸分段线性结构]\label{prop:convex}
按式\eqref{eq:q1dp} 定义的 $F_t(e)$，$t=1,\dots,145$，是其有效域（$[1200,10800]$ 与终端可达性约束之交）上的凸分段线性函数，且边际价值 $-\partial F_t(e)$ 关于 $e$ 单调非增；断点数至多逐段增加常数个。
\end{proposition}

\begin{proof}
对 $t$ 逆向归纳。$F_{145}$ 为单点指示函数，凸。设 $F_{t+1}$ 凸分段线性。记单段动作费用 $a_t(x)=p_t[L_t-V_t+\psi(x)]^+$。

(i) \emph{$a_t$ 凸分段线性}：$\psi(x)$ 是过原点的两段凸折线（斜率 $\eta<1/\eta$ 递增），$L_t-V_t+\psi(x)$ 仍凸；外层 $p_t[\cdot]^+$ 是非降凸函数与凸函数的复合，故 $a_t$ 凸，断点至多两个（$x=0$ 与 $[\cdot]^+$ 的零点）。

(ii) \emph{一步递推是下确界卷积}：换元 $u=e+x$，
\[
F_t(e)=\min_{u}\{\check a_t(e-u)+F_{t+1}(u)\}=(\check a_t\,\square\,F_{t+1})(e),\qquad \check a_t(z)=a_t(-z),
\]
再与功率限额、库存界对应的区间指示函数相加。凸分段线性函数的下确界卷积仍凸分段线性：其上境图是两上境图的 Minkowski 和，等价于把两条折线的斜率段按斜率归并后首尾相接，断点数为二者之和；与区间指示函数相加即定义域截断，保凸。故 $F_t$ 凸分段线性，断点增量不超过 $a_t$ 的段数（常数）。

(iii) \emph{边际单调}：凸函数的次微分关于自变量单调非减\cite{rockafellar}，故边际价值 $-\partial F_t$ 非增。归纳完成。
\end{proof}
